\section{Introduction}
\label{sec:introduction}
The USD/IDR rate matters to Indonesian importers, exporters, and policymakers, and geopolitical events such as conflicts and sanctions are a plausible source of short-term pressure on it. Caldara and Iacoviello build a news-based geopolitical risk index and show how risk can be measured from text~\cite{caldara2022}. Iyke et al.\ test geopolitical risk as a predictor of exchange-rate returns in 17 countries~\cite{iyke2022}. Their results do not settle the case of USD/IDR, and a risk index is a coarse summary of what the news says.

Several studies model news and prices together. Chen et al.\ select and aggregate news items with BERT representations and attention for intraday forex direction~\cite{chen2019}. Dakalbab et al.\ compare four fusion variants for forex classification, from plain concatenation to bidirectional cross-attention, and report combined-model accuracies between 0.829 and 0.852~\cite{dakalbab2025}. Their evaluation uses a 70/30 split and five-fold cross-validation. We could not verify from the accessible text whether the folds respect time order. Both use intraday data on major currency pairs, not a daily benchmark rate for the rupiah.

We ask three questions about next-fixing USD/IDR direction. (1)~Does news content add information beyond market history and the number of available articles? (2)~Which text representation works best when fusion is held fixed? (3)~With one representation fixed, do learned-query pooling, price-conditioned cross-attention, or a learned gate help?

The study makes three contributions. It builds a chronological protocol with a locked final test period, boundary purging, and training-only preprocessing. It runs a three-level ablation on one corpus of 16,723 geopolitical articles. It reports the outcome as measured: no text model beats a majority-class or price-only reference in a way that survives the uncertainty we can estimate. We treat this negative result as the main finding and describe where it applies.
